
\documentclass{article}
\usepackage{colortbl}
\usepackage{makecell}
\usepackage{multirow}
\usepackage{supertabular}

\begin{document}

\newcounter{utterance}

\centering \large Interaction Transcript for game `eqbench', experiment `eq\_bench', episode 35 with reflexion{-}Qwen3.5{-}4B.
\vspace{24pt}

{ \footnotesize  \setcounter{utterance}{1}
\setlength{\tabcolsep}{0pt}
\begin{supertabular}{c@{$\;$}|p{.15\linewidth}@{}p{.15\linewidth}p{.15\linewidth}p{.15\linewidth}p{.15\linewidth}p{.15\linewidth}}
    \# & \multicolumn{2}{c}{Player} && \multicolumn{2}{c}{Game Master} \\
    \hline

    \theutterance \stepcounter{utterance}  
    & & & \multicolumn{4}{p{0.6\linewidth}}{
        \cellcolor[rgb]{0.9,0.9,0.9}{
            \makecell[{{p{\linewidth}}}]{
                \texttt{\tiny{[P1$\langle$GM]}}
                \texttt{You have attempted this same task before and did not succeed. These are the plans you wrote after those attempts.} \\
\texttt{Trial \#1: New plan: I failed because I repeated the exact same output twice instead of providing a single, final answer. The instructions required a specific format with headings and scores, but I duplicated the block. Next time, I will analyze the dialogue to determine Scott's likely emotions based on his justification (distinguishing death from abandonment) and his offer of practical solutions, then output the scores exactly once without any repetition or extra text.} \\
\\ 
\\ 
\\ 
\texttt{Your task is to predict the likely emotional responses of a character in this dialogue:} \\
\\ 
\texttt{Deirdre: Scott, have you thought about this? Really thought about it?} \\
\texttt{Scott: Yes, mom. I have. } \\
\texttt{Deirdre: You're just going to leave me here, all alone?} \\
\texttt{Scott: That's not fair, mom. You know it's not like that.} \\
\texttt{Deirdre: Isn't it? Your father left me, and now you're doing the same.} \\
\texttt{Scott: Dad didn't leave, mom. He died. There's a difference.} \\
\texttt{Deirdre: But the end result's the same, isn't it? I'm left alone.} \\
\texttt{Scott: I'm not abandoning you. I'm just trying to live my life.} \\
\texttt{Deirdre: And what about my life, Scott? Who's going to take care of me?} \\
\texttt{Scott: You've got Mrs. Henderson next door, and I'll hire a nurse if need be.} \\
\texttt{Deirdre: I don't want a nurse, Scott. I want my son.} \\
\texttt{{[}End dialogue{]}} \\
\\ 
\texttt{At the end of this dialogue, Scott would feel...} \\
\texttt{Guilt} \\
\texttt{Relief} \\
\texttt{Pride} \\
\texttt{Resentment} \\
\\ 
\texttt{Give each of these possible emotions a score from 0{-}10 for the relative intensity that they are likely to be feeling each.} \\
\\ 
\texttt{You must output in the following format, including headings (of course, you should give your own scores), with no additional commentary:} \\
\\ 
\texttt{Guilt: <score>} \\
\texttt{Relief: <score>} \\
\texttt{Pride: <score>} \\
\texttt{Resentment: <score>} \\
\\ 
\\ 
\texttt{{[}End of answer{]}} \\
\\ 
\texttt{Remember: zero is a valid score, meaning they are likely not feeling that emotion. You must score at least one emotion > 0.} \\
\\ 
\texttt{Your answer:} \\
            }
        }
    }
     \\ \\

    \theutterance \stepcounter{utterance}  
    & \multicolumn{4}{p{0.6\linewidth}}{
        \cellcolor[rgb]{0.561,0.737,0.561}{
            \makecell[{{p{\linewidth}}}]{
                \texttt{\tiny{[P1$\rangle$GM]}}
                \texttt{Guilt: 0} \\
\texttt{Relief: 10} \\
\texttt{Pride: 5} \\
\texttt{Resentment: 2} \\
            }
        }
    }
    & & \\ \\

    \theutterance \stepcounter{utterance}  
    & & & \multicolumn{2}{p{0.3\linewidth}}{
        \cellcolor[rgb]{0.9,0.9,0.9}{
            \makecell[{{p{\linewidth}}}]{
                \texttt{\tiny{[GM$|$GM]}}
                \texttt{Guilt: 0} \\
\texttt{Relief: 10} \\
\texttt{Pride: 5} \\
\texttt{Resentment: 2} \\
            }
        }
    }
    & & \\ \\

    \theutterance \stepcounter{utterance}  
    & & & \multicolumn{2}{p{0.3\linewidth}}{
        \cellcolor[rgb]{0.9,0.9,0.9}{
            \makecell[{{p{\linewidth}}}]{
                \texttt{\tiny{[GM$|$GM]}}
                \texttt{Guilt: 8} \\
\texttt{Relief: 2} \\
\texttt{Pride: 0} \\
\texttt{Resentment: 6} \\
            }
        }
    }
    & & \\ \\

    \theutterance \stepcounter{utterance}  
    & & & \multicolumn{2}{p{0.3\linewidth}}{
        \cellcolor[rgb]{0.9,0.9,0.9}{
            \makecell[{{p{\linewidth}}}]{
                \texttt{\tiny{[GM$|$GM]}}
                \texttt{game\_result = LOSE} \\
            }
        }
    }
    & & \\ \\

    \theutterance \stepcounter{utterance}  
    & & \multicolumn{4}{p{0.6\linewidth}}{
        \cellcolor[rgb]{0.95,0.95,0.95}{
            \makecell[{{p{\linewidth}}}]{
                \texttt{\tiny{[P1$\leftrightarrow$P]}}
                \texttt{Post{-}episode Reflexion (not sent to GM):} \\
\texttt{New plan: I failed because I assigned a score of 10 to "Relief," which contradicts the narrative context. Scott is leaving his mother alone, citing her father's death and his own need to live his life, which are heavy, isolating topics. He is not relieved; he is likely feeling guilty for leaving her, resentful of her comparison to his father, and perhaps proud of his independence but not in a way that outweighs the guilt. I will re{-}analyze the dialogue to assign a high score to Guilt (due to the "abandonment" accusation) and Resentment (due to the unfair comparison), while setting Relief to 0 and adjusting Pride to reflect a complex, perhaps defensive, feeling rather than pure satisfaction. I will output the final scores exactly once in the required format.} \\
            }
        }
    }
    & \\ \\

\end{supertabular}
}

\end{document}
